\documentclass[innermargin=10mm]{tikzposter} % See Section 3
\usepackage{enumitem}
\usepackage{url}
\usepackage{wrapfig}
\geometry{paperwidth=63in,paperheight=42in}
\graphicspath{./images/}

\title{Tangible Exploration of Subsurface Data}% See Section 4.1
\institute{North Carolina State University, Raleigh, NC, United States, \texttt{akratoc@ncsu.edu}}
\author{Anna Petrasova, Brendan A. Harmon, Helena Mitasova and Jeffrey G. White}
\titlegraphic{\includegraphics{./images/logos/ncstate-type-2x2-red-rgb.eps}}

\usetheme{Autumn} % See Section 5
%  \usebackgroundstyle{Autumn}
\usetitlestyle{Default}







\definecolor{mybrown}{HTML}{58433C}
\definecolor{mybrown2}{HTML}{8B7168}
\definecolorstyle{test} {
\definecolor{colorOne}{named}{white}
\definecolor{colorTwo}{named}{mybrown}
\definecolor{colorThree}{named}{mybrown2}
}{
\colorlet{titlebgcolor}{colorOne}
\colorlet{titlefgcolor}{colorTwo}
\colorlet{framecolor}{white}
}

\usecolorstyle{test}



\usenotestyle{Corner}





\makeatletter
\renewcommand\TP@maketitle{%
    \begin{minipage}{0.2\linewidth}
       \centering
       \includegraphics[width=0.6\textwidth]{./images/logos/ncstate-type-2x2-red-rgb.eps}
    \end{minipage}%
    \hfill
   \begin{minipage}{0.6\linewidth}
        \centering
        \color{titlefgcolor}
        {\bfseries  \sc {\fontsize{85}{60}\selectfont \@title}  \par}
        \vspace*{1em}
        {\huge \@author \par}
        \vspace*{1em}
        {\LARGE \@institute}
    \end{minipage}%
    \hfill
    \begin{minipage}{0.2\linewidth}
      \centering
%        \includegraphics[width=0.37\textwidth]{grass_logo.pdf}\\
      \includegraphics[width=0.33\textwidth]{./images/logos/grass_logo.pdf}\\
      \Large {\textsf{GRASS GIS}}
   \end{minipage}
%     \begin{minipage}{0.2\linewidth}
%        \centering
%        \includegraphics[width=0.45\textwidth]{./images/logos/grass_logo.pdf}
%     \end{minipage}
}
\setlength{\TP@visibletextwidth}{\textwidth-3\TP@innermargin}
\setlength{\TP@visibletextheight}{\textheight-3\TP@innermargin}

\newcommand{\seeanimtext}{see animation online}%
\newcommand{\blocktitlesize}{\huge }
\makeatother

\setlength{\parskip}{1cm}

\begin{document}
\maketitle[width=0.7\textwidth,roundedcorners=0pt] % See Section 4.1


\begin{columns} % See Section 4.4
\column{0.2} % See Section 4.4
\block[titleleft,titleinnersep=8mm]{\blocktitlesize Introduction}{
\setlength{\parskip}{1em}
\vspace{-20pt}
\LARGE
Since traditional subsurface visualizations can be difficult to interpret
and often do not convey information in an engaging form, we have developed 
 a new tangible visualization technique which can 
 help researchers, students and the general public to
better understand and interpret volumetric subsurface data.

\begin{wrapfigure}[10]{L}{0.5\linewidth}
\vspace{-20pt}
\includegraphics[width=0.95\linewidth]{./images/setup.JPG}
\end{wrapfigure}

Our visualization technique uses the Tangible geospatial modeling and visualization system -- a system that
couples a physical 3D model with geospatial modeling through a cycle of scanning
and projection -- to study subsurface data. 
% While having a sand model of a study area,
We can manipulate the surface of a sand model representing our study area
to explore the subsurface conditions in a natural and intuitive way.
By offering a direct, tangible interface with multiple ways of interaction,
this technique is an intuitive alternative to more abstract 3D computer visualization tools.

}

\note[angle=270,radius=44cm,
width=30cm, linewidth=1cm,
roundedcorners=60]{
\LARGE

\vspace{2cm}

\textbf{Tangeoms} [1, 2]
is a collaborative modeling environment.
We can analyze the impact of terrain changes by capturing the changes on the model,
importing them into GRASS GIS [3], performing geospatial analyses
and projecting the results back on the model.
Previous applications of Tangeoms explored the impact
of terrain modifications on viewsheds, solar
energy potential, and surface-based processes, such as overland
water flow and sediment transport.

\bigskip\bigskip

\includegraphics[width=0.49\linewidth]{./images/tangeoms/state1}
\includegraphics[width=0.49\linewidth]{./images/tangeoms/state2}
}



\column{0.6}
\block[titleleft,titleinnersep=8mm]{\blocktitlesize Tangible visualization of soil ``excavation''}
{

\LARGE
\vspace{-10pt}
\begin{minipage}{0.6\linewidth}
We can look under the ground as if we were at an excavation site by digging into the sand 
with our hands and then projecting the cross-section 
of the scanned surface with the 3D raster representing soil properties 
in a real-time feedback loop. 
% you can also say 3D voxel model of soils properties
%, and see the actual data represented as a 3D raster in that particular part of the model.
In this example, we explore the 3D distribution of measured soil moisture [\%] and its change over time.
Since we are in the GRASS GIS environment, we can project additional GIS layers (such as an orthophoto, 
elevation contours, and the flow accumulation) over the model
to provide spatial context and gain more insights into soil moisture distribution 
and its relation to landscape.


\end{minipage}
\hfill
\begin{minipage}{0.37\linewidth}
\centering
\includegraphics[width=\linewidth]{./images/cross_section.png}
\end{minipage}\hfill

\bigskip\bigskip

\includegraphics[width=0.196\linewidth]{./images/excavation/2014_12_03_0042}
\includegraphics[width=0.196\linewidth]{./images/excavation/2014_12_03_0064}
\includegraphics[width=0.196\linewidth]{./images/excavation/2014_12_03_0065}
\includegraphics[width=0.196\linewidth]{./images/excavation/2014_12_03_0070}
\includegraphics[width=0.196\linewidth]{./images/excavation/2014_12_03_0073}
\raisebox{0pt}[0pt][0pt]{\hspace{-3cm}\colorbox{white}{\includegraphics[width=3cm]{images/qrcodes/youtube1}}}%
\newlength{\myl}%
\settowidth{\myl}{\tiny\seeanimtext}%
\raisebox{-1ex}[0pt][0pt]{\hspace{-1.8\myl}\colorbox{white}{\small\seeanimtext}}

}


\block[titleleft,titleinnersep=8mm]{\blocktitlesize Visualizing vertical profiles}
{
\LARGE
\vspace{-10pt}
\begin{minipage}{0.745\linewidth}
\centering

% \includegraphics[width=0.25\linewidth]{./images/slices/2014_12_03_0096_cut}\hfill
\includegraphics[width=0.32\linewidth]{./images/slices/2014_12_03_0098_cut}
\includegraphics[width=0.32\linewidth]{./images/slices/2014_12_03_0101_cut}
\includegraphics[width=0.32\linewidth]{./images/slices/2014_12_03_0111_cut}

\bigskip\bigskip

\includegraphics[width=0.32\linewidth]{./images/slices/2014_12_03_0123_cut}
\includegraphics[width=0.32\linewidth]{./images/slices/2014_12_03_0129_cut}
\includegraphics[width=0.32\linewidth]{./images/slices/2014_12_03_0130_cut}
\end{minipage}
\begin{minipage}{0.25\linewidth}
By placing two markers, we specify the location and length of a vertical subsurface profile
which is then displayed in the legend. This technique is optionally combined with the ``excavation'' visualization.
Furthermore, we can display other soil properties
(such as moisture together with the distribution of clay, silt and sand particles)
to discover relationships between different subsurface properties.
\end{minipage}
}


\block[titleleft,titleinnersep=8mm]{\blocktitlesize ``Core sample'' visualization}
{
\vspace{-10pt}
\LARGE
\begin{minipage}{0.28\linewidth}
We can query the 3D raster vertically at sites specified by placing markers, 
similar to how we used markers to examine soil cores at given locations. 
\end{minipage}
\begin{minipage}{0.72\linewidth}
\centering

\includegraphics[width=0.24\linewidth]{./images/cores/2014_12_03_0141}
\includegraphics[width=0.24\linewidth]{./images/cores/2014_12_03_0144}
\includegraphics[width=0.24\linewidth]{./images/cores/2014_12_03_0151}
\includegraphics[width=0.24\linewidth]{./images/cores/2014_12_03_0154}
\raisebox{0pt}[0pt][0pt]{\hspace{-3cm}\colorbox{white}{\includegraphics[width=3cm]{images/qrcodes/youtube2}}}%
\settowidth{\myl}{\small\seeanimtext}%
\raisebox{-1ex}[0pt][0pt]{\hspace{-\myl}\colorbox{white}{\small\seeanimtext}}
\end{minipage}
}


\block[titleleft,titleinnersep=8mm]{\blocktitlesize Time series visualization}
{
\vspace{-10pt}
\LARGE
\begin{minipage}{0.28\linewidth}
It is possible to animate measured soil moisture time series
using any combination of our visualization techniques.
The date of measurement is displayed together with the data and the legend.
\end{minipage}
\begin{minipage}{0.72\linewidth}
\centering

\includegraphics[width=0.24\linewidth]{./images/timeseries/state2}
\includegraphics[width=0.24\linewidth]{./images/timeseries/state3}
\includegraphics[width=0.24\linewidth]{./images/timeseries/state4}
\includegraphics[width=0.24\linewidth]{./images/timeseries/state5}
\raisebox{0pt}[0pt][0pt]{\hspace{-3cm}\colorbox{white}{\includegraphics[width=3cm]{images/qrcodes/youtube3}}}%
\settowidth{\myl}{\small\seeanimtext}%
\raisebox{-1ex}[0pt][0pt]{\hspace{-\myl}\colorbox{white}{\small \seeanimtext}}

\end{minipage}
}






\column{0.2}
\block[titleleft,titleinnersep=8mm]{\blocktitlesize Soil data}{
\begin{wrapfigure}[5]{R}{0.5\linewidth}
\vspace{-20pt}\hfill
\includegraphics[width=0.95\linewidth]{./images/soil_map.JPG}
\hfill
\end{wrapfigure}
\setlength{\parskip}{0.5cm}
\LARGE
The study site is a 12~ha field located at the Lower
Coastal Plain Tobacco Research Station, Kinston, North Carolina.
Measured soil physical properties include texture, bulk density, porosity
and volumetric soil moisture (measured by TDR probes).
Properties were measured at 60 locations and five depths (8, 23, 38, 53, and 68 cm);
soil moisture measurements started in 2002 and ended in 2005 with irregular sampling intervals [4].


We imported measured data in GRASS GIS as 3D points and an interpolated 3D raster (volume)
using the regularized spline with tension function implemented in the GRASS GIS module v.vol.rst [5].
The depth was exaggerated by a factor of 100.
}

\block[titleleft,titleinnersep=8mm]{\blocktitlesize Additional resources}{
% qrcode -t EPS -o testing_doc.eps "www.youtube.com/watch?v=QBPzYXiL3TY"

\begin{description}
 \item[Related videos] ~\\
 \url{www.youtube.com/watch?v=QBPzYXiL3TY}\\
 \url{www.youtube.com/watch?v=__fWppH9aqQ}\\
 \url{www.youtube.com/watch?v=dy6U1jnYkl0}
 
 \item[NCSU OSGeo Research and Education Lab]~\\
 \url{geospatial.ncsu.edu/osgeorel/tangeoms.html}
 
 \item[GRASS GIS 3D raster data processing]~\\{\fontsize{29}{20}\selectfont
 \url{grass.osgeo.org/grass70/manuals/raster3dintro.html}}
 
\end{description}

}
\block[titleleft]{\blocktitlesize References}{
\fontsize{27}{37.} \selectfont
\begin{enumerate}[label={[\arabic*]}]
 \item Petrasova, A., Harmon, B. A., Petras, V., Mitasova, H., 2014.
\emph{GIS-based environmental modeling with tangible interaction and dynamic visualization.}
 In: Ames, D.P., Quinn, N.W.T., Rizzoli, A.E. (Eds.),
 Proceedings of the 7th International Congress on Environmental Modelling and Software,
 June 15--19, San Diego, California, USA.

\item Tateosian, L., Mitasova, H., Harmon, B. A., Fogleman, B., Weaver, K., Harmon, R. S., 2010.
\emph{TanGeoMS: tangible geospatial modeling system.}
IEEE Transactions on Visualization and Computer Graphics, 16(6), 1605--12.

\item Neteler, M., Bowman, M. H., Landa, M., Metz, M., 2012.
\emph{GRASS GIS: A multi-purpose open source GIS.}
Environmental Modelling \& Software, 31(0), 124--130.

\item Duffera, M., White, J. G., Weisz, R., 2007.
\emph{Spatial variability of Southeastern U.S. Coastal Plain soil physical properties:
Implications for site-specific management.}
Geoderma, 137(3-4), 327--339.

\vspace{5pt}

\item Hofierka, J., Parajka, J., Mitasova, H., Mitas, L., 2002.
\emph{Multivariate interpolation of precipitation using regularized spline with tension.}
Transactions in GIS, 6(2), 135--150.
\end{enumerate}
}
\end{columns}
\end{document}